% ===================================================================
% Background and Related Work — revised to focus on literature that
% directly motivates the thesis contribution.
% ===================================================================

Predictive multiplicity lies at the intersection of three lines of work: the study of
Rashomon sets, the design of metrics that quantify disagreement among near-optimal
models, and the broader question of what such disagreement implies for model selection,
fairness, and reliability. The central observation is that a single deployed model is often
only one choice among several similarly well-performing alternatives. Even when a
machine-learning workflow ultimately returns one final model, other choices of model
family, hyperparameters, or random seed may lead to models with nearly identical
predictive performance but different outputs for the same individual. This chapter
reviews the literature most directly connected to that observation and uses it to motivate
the perspective taken in this thesis: predictive multiplicity should be studied not only as
a global property of a dataset or model class, but also as a localized phenomenon in
feature space.

% -------------------------------------------------------------------
\subsection{Rashomon Sets and Predictive Multiplicity}
% -------------------------------------------------------------------

The modern formulation of multiplicity is rooted in the \emph{Rashomon effect},
introduced by \cite{breiman_statistical_2001}. The Rashomon effect describes the situation in which many distinct models achieve
similarly good predictive performance, while relying on different functional forms,
variables, or relationships in the data. In other words, there may be no single
uniquely best explanation of the data-generating relationship. This perspective is also
used by \cite{fisher2019all}, who argue that conclusions drawn from one fitted model
can be incomplete when other well-performing models rely on different variables or
relationships.

Formally, let $\mathcal{F}$ be a hypothesis class, let $L(f)$ be the loss of a model
$f \in \mathcal{F}$, and let $L^* = \min_{f \in \mathcal{F}} L(f)$ be the best loss in this
class. The corresponding \emph{Rashomon set} is
\[
  \mathcal{R}_\varepsilon
  = \bigl\{f \in \mathcal{F} : L(f) \leq L^{*} + \varepsilon\bigr\}.
\]
It contains all models whose loss is within a tolerance $\varepsilon$ of the best model.
This shifts attention from a single selected model to a set of near-optimal alternatives.

Recent Rashomon-set research has emphasized that the size and structure of
$\mathcal{R}_\varepsilon$ can have substantial practical consequences.
\cite{semenova_existence_2022} formalize the \emph{Rashomon ratio}, i.e., the volume of
$\mathcal{R}_\varepsilon$ relative to the hypothesis space, and argue that large Rashomon
sets help explain why simpler yet accurate models may exist.
\cite{rudin_amazing_2024} broaden this perspective further: when many good models are
available, model development should not be viewed as a search for a uniquely correct
solution, but as a search over a set of viable models that may differ in interpretability,
fairness, robustness, and other secondary properties.

Within this broader setting, \cite{marx_predictive_2020} introduce
\emph{predictive multiplicity} as the phenomenon that near-optimal models can assign
conflicting predictions to the same observation.
\cite{black_model_2022} generalize this concern under the broader notion of
\emph{model multiplicity}, highlighting that arbitrariness in model choice can create
both opportunities and risks: multiplicity may enable deliberate optimization for fairness
or interpretability, but it may also make individual outcomes depend on choices that are
invisible to affected stakeholders.

Predictive multiplicity should also be distinguished from classical notions of
uncertainty. Aleatoric uncertainty reflects inherent noise in the data, while epistemic
uncertainty captures limited knowledge about the data-generating process. In contrast,
predictive multiplicity arises from the existence of multiple equally accurate models
that produce different predictions for the same observation. The resulting disagreement
is therefore not due only to noise or lack of information, but to the structure of the
model space itself.

This thesis studies precisely that predictive facet of multiplicity: how much near-optimal
models disagree on individual observations, where that disagreement is concentrated, and
which modeling choices drive it.

% -------------------------------------------------------------------
\subsection{Metrics for Predictive Multiplicity}
\label{sec:metrics-review}
% -------------------------------------------------------------------

The first generation of predictive-multiplicity metrics operates in \emph{label space}.
\cite{marx_predictive_2020} propose \emph{ambiguity}, which measures how often a sample
can receive conflicting labels within the Rashomon set, and \emph{discrepancy}, which
quantifies the maximal extent of such disagreement. These measures made predictive multiplicity operational, but they reduce model outputs to binary decisions based on a fixed threshold (e.g., 0.5), thereby ignoring differences in predicted probabilities when models agree on the final label.

A second line of work extends multiplicity to \emph{probability space}.
\cite{watson-daniels_predictive_2023} develop a framework for predictive multiplicity in
probabilistic classification, introducing score-based measures that capture the variation
of risk estimates across near-optimal models rather than only label flips. Their analysis
also connects multiplicity to dataset characteristics such as separability, outliers, and
majority--minority structure. This is especially relevant for this thesis because the
main response variable is not hard disagreement alone, but the full spread of predicted
probabilities across the Rashomon set.

Closely related, \cite{hsu_rashomon_nodate} propose \emph{Rashomon Capacity}, an
information-theoretic measure of predictive multiplicity in probabilistic classification.
The measure is defined at the level of an individual observation and captures how much
the possible predicted scores for that observation vary across models in the Rashomon set.
\cite{hsu_rashomongb_2024} extend this line in the context of gradient boosting,
introducing \emph{RashomonGB} as a practical method for inspecting multiplicity in boosted
models and discussing score-based notions such as the \emph{Viable Prediction Range} for
individual observations.

At the per-observation level, one particularly useful hard-disagreement measure is the
\emph{conflict ratio}, which records the fraction of models on the minority side of a
decision threshold. Recent work by \cite{frohnapfel_using_2026} argues that such
individual-level multiplicity measures are not only methodologically useful but also
policy-relevant, because they can help operationalize per-person performance reporting
requirements in high-stakes settings under the EU AI Act.

Taken together, the literature now offers several ways to quantify multiplicity globally
and per instance, in both label space and probability space. What remains much less
developed is how to move from pointwise disagreement scores to an analysis of whether
disagreement forms coherent \emph{regions} in feature space.


\subsection{Localizing Instability in Feature Space}
\label{sec:spatial-review}
% -------------------------------------------------------------------

The metrics discussed above quantify predictive multiplicity either at the dataset level
or for individual observations. A remaining question is whether observations with high
prediction disagreement are isolated cases, or whether they form coherent regions in
feature space. This thesis approaches that question using tools from spatial statistics.

\cite{moran1950notes} introduced Moran's $I$ as a measure of global spatial
autocorrelation. In this thesis, it is used to test whether observations with similar
levels of prediction instability tend to be close to each other in feature space.
\cite{anselin1995local} developed \emph{Local Indicators of Spatial Association} (LISA),
which identify where such local clustering occurs. In particular, LISA can identify
\emph{High--High} regions, where observations with high values are surrounded by
neighbors that also have high values, and \emph{Low--Low} regions, where low values
cluster together.

In their original setting, these tools operate on geographic data with an explicit spatial
adjacency structure. In this thesis, the same logic is transferred to
\emph{feature space}: observations become nodes in a $k$-nearest-neighbor graph, and
observation-wise multiplicity scores become the variable whose spatial organization is
tested. Conceptually, this turns predictive multiplicity from a collection of isolated
pointwise values into a structured signal that can be analyzed for clustering, significance,
and local hotspot formation.

This step is motivated by observations already present in the literature on predictive
multiplicity. \cite{watson-daniels_predictive_2023} show that prediction disagreement
across near-optimal probabilistic classifiers is associated with dataset characteristics such
as outliers and separability. This suggests that predictive multiplicity may not be
distributed uniformly across the input space. Likewise, per-instance feasible-range
methods such as those discussed by \cite{hsu_rashomongb_2024} show that the range of
possible predictions can vary substantially from one observation to another. These works
therefore motivate an observation-level view of predictive multiplicity. However, they do
not directly analyze whether observations with high disagreement form statistically
significant, spatially contiguous regions in feature space.

That question matters for interpretation. Once disagreement is shown to cluster, the
relevant object is no longer only an unstable individual prediction, but an unstable
\emph{subpopulation}. This shifts the focus from isolated observations to coherent
regions of the feature space, enabling a more structured analysis of where and how
predictive multiplicity arises.


% -------------------------------------------------------------------
\subsection{Drivers and Consequences of Predictive Multiplicity}
\label{sec:drivers-review}
% -------------------------------------------------------------------

After asking how predictive multiplicity can be measured and localized, a further question
is why it arises and why it matters. One source is model selection itself. If several models
perform similarly well, then choosing one final model is not only a technical step, but can
affect individual predictions. \cite{black_model_2022} and \cite{rudin_amazing_2024}
emphasize this point: when many good models are available, model choice should be seen
as a substantive design decision, because different acceptable models may differ in other
properties such as interpretability, fairness, or robustness.

More recently, attention has turned to concrete modeling choices that generate
multiplicity within a model family. Most directly relevant for this thesis is
\cite{cavus_role_2025}, who show across benchmark datasets that hyperparameter tuning
can materially increase prediction discrepancy even when performance improves.
Their results suggest that multiplicity is not driven only by broad differences between
families such as linear models versus tree ensembles, but also by movement through
hyperparameter space within a fixed family.

Recent work also highlights that prediction variance across models can lead to
arbitrary decisions for individual observations \citep{cooper2024arbitrariness}.

The implications of such disagreement have also been studied from a fairness and
governance perspective. \cite{sokol_cross-model_2024} formulate
\emph{cross-model fairness}, emphasizing that when different equally performing models
treat the same person differently, fairness concerns arise not only from the deployed
predictor itself but also from the unresolved choice among competing predictors.
Similarly, \cite{frohnapfel_using_2026} connect predictive multiplicity to legal and
reporting obligations in high-risk decision support, arguing that disagreement across
near-optimal models can be interpreted as a form of person-level arbitrariness that should
be made visible to deployers.

These works clarify why multiplicity matters and which modeling choices may create it.
However, they mostly operate at the level of whole datasets, whole model classes, or
individual cases taken one at a time. They do not yet ask whether family effects and
hyperparameter effects become especially strong in \emph{specific parts} of feature space.


% -------------------------------------------------------------------
\subsection{Positioning of This Thesis}
% -------------------------------------------------------------------

The existing literature provides concepts and metrics for studying predictive
multiplicity, but it mostly treats disagreement either as a dataset-level quantity or as a
property of individual observations. Less attention has been given to whether these
observation-level differences form coherent regions in feature space.

This thesis addresses this gap by treating prediction instability as a spatial signal. It uses
observation-wise predictive variance to identify where near-optimal models disagree, and
then analyzes whether high-disagreement observations cluster locally. The goal is therefore not only to measure how much predictive
multiplicity exists, but also to locate where it occurs, assess how stable these regions
are, and examine which modeling choices are associated with them.

In this way, the thesis connects predictive multiplicity with a regional view of model
reliability: some parts of the feature space may be more sensitive to model choice than
others, even when all considered models perform similarly well overall.